% ---------------------------------------------------------------------
\section{How to Read This Guide}
% ---------------------------------------------------------------------

HyperPod turns a collection of ordinary laptops into one shared computer that
runs a large language model (LLM). This guide documents the \emph{idea}, the
\emph{evidence} behind it, and the exact \emph{plan} to build it for the Meta
Global AI Developer Hackathon (Indie track). It is written so a newcomer can
follow it, while remaining strictly honest about what is proven and what is not.

\subsection{The evidence tags}
Next to important claims you will see one of three tags. They answer the
question ``how much should I trust this?'':

\begin{itemize}[leftmargin=1.4em,nosep]
  \item \VER\ --- a \textbf{verified fact} taken from an official repository,
  research paper, or the hackathon's own rules. These can be stated to a judge.
  \item \HYP\ --- a \textbf{reasoned hypothesis}: a sensible inference that has
  not yet been proven by a test or a source.
  \item \UNK\ --- \textbf{unknown}: not enough evidence yet; must be checked
  before relying on it.
\end{itemize}

\subsection{Keep vs drop}
Design decisions are marked \Keep\ (part of the plan) or \Drop\ (deliberately
excluded, with a reason). Being explicit about what we are \emph{not} building
is as important as what we are.

\subsection{The one-paragraph mental model}
A single LLM is split into layers. Each laptop (a ``worker'') loads a slice of
those layers. To answer a prompt, data flows worker~1 → worker~2 → worker~3 like
an assembly line, and words come out the end. A central ``controller'' hands out
the work, listens for heartbeats, and --- the key part --- when a worker goes
silent, it reassigns that worker's slice to a healthy machine and re-runs the
request so the answer continues. The whole point is that the laptops are
\emph{unreliable by nature} (they are people's real machines on Wi-Fi), and
HyperPod makes the cluster behave reliably anyway.

\subsection{The honest novelty statement}
We do not claim to have invented load-aware failure recovery for LLMs. That
already exists in datacenter research (\S\ref{sec:landscape}). HyperPod's
contribution, stated safely, is: \emph{``Fault tolerance like this exists in
datacenters. It has never run on a pile of cheap student laptops over Wi-Fi with
free tools. We made that work --- and here, watch, we kill one.''} This framing
is defensible because every part of it is backed by a source in
\S\ref{sec:landscape}.
